% Hình kiến trúc front_qconv4x3_proj10 (level32 + encoder quantum qconv4x3), 3 trang:
%   trang 1: tổng quan luồng dữ liệu; trang 2: một tầng encoder qconv; trang 3: một cây tensor network.
% Biên dịch: pdflatex hinh_kien_truc_qconv4x3.tex (Overleaf: compiler pdfLaTeX). Cần gói vntex (có sẵn trong TeX Live).
% Giải thích chữ: KIEN_TRUC_front_qconv4x3_proj10.txt
\documentclass[tikz,border=10pt]{standalone}
\usepackage[T5]{fontenc}
\usepackage[utf8]{inputenc}
\usepackage[vietnamese]{babel}
\usepackage{amsmath,amssymb}
\usetikzlibrary{positioning,arrows.meta,calc,fit,backgrounds,decorations.pathreplacing}

\tikzset{
  >={Stealth[length=2.2mm]},
  qbox/.style={draw=blue!60!black, fill=blue!9, rounded corners=2pt, align=center, font=\small, inner sep=4pt},
  cbox/.style={draw=orange!75!black, fill=orange!13, rounded corners=2pt, align=center, font=\small, inner sep=4pt},
  dbox/.style={draw=black!55, fill=black!4, rounded corners=2pt, align=center, font=\small, inner sep=4pt},
  note/.style={align=left, font=\footnotesize, text=black!75},
  flow/.style={->, thick, black!70},
  gate/.style={draw=blue!60!black, fill=blue!12, minimum width=7mm, minimum height=5mm, inner sep=1pt, font=\scriptsize},
  meas/.style={draw=blue!60!black, fill=white, minimum width=14mm, minimum height=5mm, inner sep=1pt, font=\scriptsize},
}

% Lưới 4x4 (hàng 0 ở trên) với đường quét qua tâm các ô: #1 = x, #2 = tên, #3 = danh sách tọa độ (cột, hàng).
\newcommand{\scangrid}[3]{%
  \begin{scope}[shift={(#1,-6.8)}, x=0.55cm, y=-0.55cm]
    \foreach \i in {0,...,4} {
      \draw[black!30] (\i,0) -- (\i,4);
      \draw[black!30] (0,\i) -- (4,\i);
    }
    \draw[->, thick, blue!60!black] plot coordinates {#3};
    \fill[red!70!black] (0.5,0.5) circle (2.2pt);
    \node[font=\small] at (2,4.6) {#2};
  \end{scope}%
}

\begin{document}

% =====================================================================================================
% TRANG 1: TỔNG QUAN
% =====================================================================================================
\begin{tikzpicture}
  \node[font=\large\bfseries, anchor=west] at (0,4.3) {front\_qconv4x3\_proj10: luồng dữ liệu (MNIST 16$\times$16, 17.264 tham số)};

  % ảnh đầu vào
  \fill[black!4] (0,-1.5) rectangle (3,1.5);
  \draw[black!35, step=0.1875] (0,-1.5) grid (3,1.5);
  \draw[black!70] (0,-1.5) rectangle (3,1.5);
  \node[dbox, anchor=north, text width=3.2cm] (inlab) at (1.5,-1.8)
    {Ảnh 28$\times$28 $\to$ 16$\times$16\\(aug khi train)\\256 pixel, $\theta_p$ = giá trị,\\$\phi_p = 0$};
  \coordinate (img) at (3,0);

  % encoder
  \node[qbox, text width=3.1cm] (s1) at (6.2,1.4) {Tầng 1\\cửa sổ 3$\times$3};
  \node[qbox, text width=3.1cm] (s2) at (6.2,0)   {Tầng 2\\cửa sổ 5$\times$5};
  \node[qbox, text width=3.1cm] (s3) at (6.2,-1.4) {Tầng 3\\cửa sổ 7$\times$7};
  \draw[flow] (s1) -- (s2);
  \draw[flow] (s2) -- (s3);
  \begin{scope}[on background layer]
    \node[draw=blue!60!black, dashed, rounded corners, fit=(s1)(s3), inner sep=7pt, label={[font=\small\bfseries, align=center]above:Encoder quantum qconv4x3\\(492 tham số)}] (enc) {};
  \end{scope}
  \node[note, below=2mm of enc, text width=4.2cm]
    {Mỗi pixel: mạch 9 qubit trên ô 3$\times$3,\\4 lớp, đo $X,Y,Z$ (27 số),\\Linear(27,2) cộng dư vào $(\theta_p,\phi_p)$.\\Xem trang 2.};
  \draw[flow] (img) -- (enc.west);

  % mã hóa
  \node[qbox, text width=3.6cm] (code) at (11.2,0)
    {Mã hóa 256 qubit\\[2pt]$|q_p\rangle = \cos\tfrac{\pi\theta_p}{2}|0\rangle$\\$+\,e^{i\pi\phi_p}\sin\tfrac{\pi\theta_p}{2}|1\rangle$};
  \draw[flow] (enc.east) -- node[above, font=\footnotesize]{$\theta,\phi$} node[below, font=\footnotesize]{(256, 256)} (code.west);

  % 3 cây
  \node[qbox, text width=3.3cm] (t1) at (16.2,2.3)  {Cây 1: quét theo hàng\\256 $\to$ 8 qubit};
  \node[qbox, text width=3.3cm] (t2) at (16.2,0)    {Cây 2: quét theo cột\\256 $\to$ 8 qubit};
  \node[qbox, text width=3.3cm] (t3) at (16.2,-2.3) {Cây 3: quét Morton\\256 $\to$ 8 qubit};
  \foreach \t in {t1,t2,t3} \draw[flow] (code.east) -- (\t.west);
  \node[note, below=2mm of t3, text width=3.6cm] {Mỗi cây 1.608 tham số (xem trang 3).\\Cùng 256 qubit, chỉ đổi thứ tự lá.};

  % projection mỗi cây
  \foreach \t/\i in {t1/1,t2/2,t3/3} {
    \node[cbox, text width=2.5cm, right=1.5cm of \t] (p\i) {Linear(372,10)\\+ LN + GELU};
    \draw[flow] (\t.east) -- node[above, font=\footnotesize]{372} (p\i.west);
  }
  \node[note, above=1mm of p1] {372 = 276 (8 qubit cuối) + 96 (tầng 32)};

  % gộp
  \node[cbox, text width=2.9cm] (fus) at (25.4,0) {Nối 3$\times$10 = 30\\Linear(30,16)\\+ LN + GELU\\Linear(16,10)};
  \foreach \i in {1,2,3} \draw[flow] (p\i.east) -- (fus.west);
  \node[dbox, right=1cm of fus] (out) {10 logit\\cross-entropy};
  \draw[flow] (fus) -- (out);

  % chú giải
  \node[qbox, minimum width=6mm, minimum height=4mm] (lq) at (0.4,-5.2) {};
  \node[anchor=west, font=\small] at (lq.east) {quantum (mạch mô phỏng chính xác)};
  \node[cbox, minimum width=6mm, minimum height=4mm] (lc) at (7.6,-5.2) {};
  \node[anchor=west, font=\small] at (lc.east) {cổ điển};
  \node[anchor=west, font=\small, align=left] at (0,-6.2)
    {Tham số: quantum 5.148 (29,8\%) = góc mạch encoder 324 + 3 cây 4.824;\quad
     cổ điển 12.116 = readout encoder 168 + head 11.948.\\
     Kết quả (val, seed 0, dropout 0, giữa run ep 32): MAX 99,03\% = CK 99,03\% (ep 29).};
\end{tikzpicture}

% =====================================================================================================
% TRANG 2: MỘT TẦNG ENCODER qconv (lặp 3 lần, tham số riêng mỗi tầng)
% =====================================================================================================
\begin{tikzpicture}
  \node[font=\large\bfseries, anchor=west] at (0,1.6) {Một tầng encoder qconv (qconv4x3 có 3 tầng như vậy, mỗi tầng 164 tham số)};

  % ô 5x5 quanh pixel p, cửa sổ 3x3 tô màu
  \begin{scope}[shift={(0,-4.6)}]
    \fill[blue!12] (1,1) rectangle (4,4);
    \draw[black!40] (0,0) grid (5,5);
    \draw[blue!60!black, very thick] (1,1) rectangle (4,4);
    \foreach \r in {0,1,2} \foreach \c in {0,1,2} {
      \pgfmathtruncatemacro{\k}{3*\r+\c}
      \node[font=\small] at (1.5+\c, 3.5-\r) {$q_{\k}$};
    }
    \draw[red!70!black, thick] (2,2) rectangle (3,3);
    \node[note, below=1mm, text width=5cm, align=center] at (2.5,0)
      {Cửa sổ 3$\times$3 quanh pixel $p$ (= $q_4$).\\Ngoài ảnh: qubit $|0\rangle$.\\Cùng mạch cho cả 256 pixel (như tích chập).};
  \end{scope}

  % đồ thị CZ 3x3
  \begin{scope}[shift={(0.9,-9.4)}]
    \foreach \r in {0,1,2} \foreach \c in {0,1,2} \node[circle, draw=blue!60!black, fill=blue!12, inner sep=1.5pt, font=\tiny] (n\r\c) at (1.6*\c, -1.2*\r) {};
    \foreach \r in {0,1,2} { \draw[thick] (n\r0) -- (n\r1) -- (n\r2); }
    \foreach \c in {0,1,2} { \draw[thick] (n0\c) -- (n1\c) -- (n2\c); }
    \node[note, align=center] at (1.6,-3.1) {CZ trên 12 cạnh lưới 3$\times$3\\(ngang, dọc; không chéo)};
  \end{scope}

  % mạch 9 qubit
  \def\xs{7.0}
  \def\xe{19.6}
  \foreach \k in {0,...,8} {
    \pgfmathsetmacro{\y}{-0.75*\k}
    \draw[black!70] (\xs,\y) -- (\xe,\y);
    \node[anchor=east, font=\small] at (\xs,\y) {$q_{\k}$};
    \node[gate] at (8.0,\y) {$E_{\k}$};
    \node[gate] at (9.6,\y) {$R$};
    \node[gate] at (13.0,\y) {$R$};
    \node[meas] at (16.7,\y) {$\langle X\rangle\langle Y\rangle\langle Z\rangle$};
  }
  \foreach \x in {10.9,14.3} {
    \node[draw=blue!60!black, fill=blue!20, minimum width=7mm, minimum height=6.5cm, font=\scriptsize] at (\x,-3.0) {CZ};
  }
  \foreach \k in {0,...,8} { \node[fill=white, inner sep=1pt] at (11.95,-0.75*\k) {$\cdots$}; }

  \draw[decorate, decoration={brace, amplitude=5pt}, thick] (9.2,0.45) -- node[above=5pt, font=\footnotesize]{lớp 1} (11.3,0.45);
  \draw[decorate, decoration={brace, amplitude=5pt}, thick] (12.6,0.45) -- node[above=5pt, font=\footnotesize]{lớp 4} (14.7,0.45);
  \node[font=\footnotesize] at (11.95,0.95) {lớp 2, 3};

  % readout
  \draw[decorate, decoration={brace, amplitude=6pt}, thick] (17.8,0.3) -- (17.8,-6.3);
  \node[font=\footnotesize, align=center] at (18.7,-3.0) {27\\số};
  \node[cbox, text width=2.2cm] (lin) at (21.3,-3.0) {Linear(27,2)\\(56 tham số,\\khởi tạo 0)};
  \draw[flow] (19.3,-3.0) -- (lin.west);
  \node[dbox, text width=2.6cm, right=0.9cm of lin] (upd) {$\theta_p \leftarrow \theta_p + \Delta\theta$\\$\phi_p \leftarrow \phi_p + \Delta\phi$};
  \draw[flow] (lin) -- node[above, font=\footnotesize]{$\Delta\theta,\Delta\phi$} (upd);

  % chú thích cổng
  \node[note, anchor=north west, text width=15cm] at (7.0,-7.2)
    {$E_k = R_Z(\pi\phi_k)\,R_Y(\pi\theta_k)$: mã hóa pixel thứ $k$ của cửa sổ, $|0\rangle \to \cos\tfrac{\pi\theta_k}{2}|0\rangle + e^{i\pi\phi_k}\sin\tfrac{\pi\theta_k}{2}|1\rangle$.\\[2pt]
     $R = R_Z(a)\,R_Y(b)\,R_Z(c)$, góc riêng cho từng qubit và từng lớp: $4 \times 9 \times 3 = 108$ góc học được mỗi tầng.\\[2pt]
     Mô phỏng statevector chính xác $2^9 = 512$ biên độ; đo kỳ vọng $X, Y, Z$ của 9 qubit.\\[2pt]
     Đo xong rồi mã hóa lại ở tầng sau: không có vướng víu nối giữa các cửa sổ hay giữa các tầng.\\[2pt]
     Readout khởi tạo 0 nên lúc đầu encoder là ánh xạ đồng nhất (model khởi đầu giống level32).\\[2pt]
     3 tầng nối tiếp: cửa sổ nhìn được tăng 3$\times$3 $\to$ 5$\times$5 $\to$ 7$\times$7.\quad Mỗi tầng: 108 + 56 = 164; 3 tầng: 492.};
\end{tikzpicture}

% =====================================================================================================
% TRANG 3: MỘT CÂY TENSOR NETWORK (mỗi cây 1.608 tham số; 3 cây 4.824)
% =====================================================================================================
\begin{tikzpicture}
  \node[font=\large\bfseries, anchor=west] at (-6,1.4) {Một cây tensor network (level32): 256 $\to$ 8 qubit, 1.608 tham số};

  % các tầng của cây (thanh hẹp dần)
  \foreach \n/\w/\y in {256/12/0, 128/10/-1.4, 64/8/-2.8, 32/6/-4.2, 16/4/-5.6, 8/2.4/-7.0} {
    \node[qbox, minimum width=\w cm, minimum height=6mm, inner sep=1pt] (L\n) at (0,\y) {\n\ qubit};
  }
  \foreach \a/\b/\i in {256/128/1, 128/64/2, 64/32/3, 32/16/4, 16/8/5} {
    \draw[flow] (L\a) -- node[left=2pt, font=\footnotesize]{conv-pool \i} (L\b);
  }
  \node[note, anchor=west] at (6.3,0) {Lá: 256 qubit đã mã hóa, theo thứ tự quét của cây};

  % đo tầng 32
  \node[cbox, text width=3.4cm] (side) at (6.4,-4.2) {Đo $\langle X\rangle,\langle Y\rangle,\langle Z\rangle$\\từng qubit: 96 số};
  \draw[flow, red!70!black] (L32.east) -- (side.west);

  % sau 8 qubit
  \node[qbox, text width=6cm, below=7mm of L8] (re) {Re-upload: mỗi qubit cuối đại diện 32 pixel liên tiếp;\\$R_Y(s_q \cdot \bar\theta_q)$, $s_q$ học được (8 tham số, khởi tạo 0)};
  \node[qbox, text width=6cm, below=5mm of re] (se) {Strongly-entangling 4 lớp trên 8 qubit\\(ma trận mật độ đủ $256\times256$; 1.280 tham số)\\Chỉ ở đây mới có vướng víu giữa các qubit};
  \node[qbox, text width=6cm, below=5mm of se] (pm) {Đo Pauli: $X,Y,Z$ của 8 qubit (24)\\+ 9 tương quan $XX,\dots,ZZ$ của 28 cặp (252) = 276};
  \draw[flow] (L8) -- (re);
  \draw[flow] (re) -- (se);
  \draw[flow] (se) -- (pm);
  \node[dbox, text width=3.4cm] (cat) at (6.4,-11.9) {Nối 276 + 96\\= 372 số $\to$ head};
  \draw[flow] (pm.east) -- (cat.west);
  \draw[flow, red!70!black] (side.south) -- (cat.north);

  % chi tiết một khối conv-pool
  \begin{scope}[shift={(11.5,-1.2)}]
    \node[font=\small\bfseries, anchor=west] at (0,0.9) {Một khối conv-pool (trên mỗi cặp qubit kề nhau)};
    \draw[black!70] (0,0) -- (6.2,0);
    \draw[black!70] (0,-1) -- (4.1,-1);
    \node[anchor=east, font=\small] at (0,0) {$\rho_{2j}$};
    \node[anchor=east, font=\small] at (0,-1) {$\rho_{2j+1}$};
    \node[draw=blue!60!black, fill=blue!15, minimum width=9mm, minimum height=17mm] at (1.4,-0.5) {$U_c$};
    \node[draw=blue!60!black, fill=blue!15, minimum width=9mm, minimum height=17mm] at (2.9,-0.5) {$U_p$};
    \node[draw, fill=white, font=\scriptsize, inner sep=2pt] at (4.3,-1) {Tr};
    \node[anchor=west, font=\small] at (6.2,0) {1 qubit (hỗn hợp)};
    \node[note, anchor=north west, text width=7.6cm] at (-0.6,-1.6)
      {$U = \exp(iH)$, $H$ Hermite từ ma trận phức $4\times4$ học được (32 số).\\
       Mỗi tầng một cặp $(U_c, U_p)$ = 64 tham số, dùng chung cho mọi cặp; 5 tầng = 320.\\
       Lấy vết bỏ qubit thứ hai ngay, nên các qubit trong cùng tầng luôn độc lập:\\
       vector Bloch $(X,Y,Z)$ từng qubit là toàn bộ thông tin của tầng.};
  \end{scope}

  % ba thứ tự quét (ví dụ ảnh 4x4)
  \node[font=\small\bfseries, anchor=west] at (11.5,-6.3) {Ba thứ tự quét (minh họa trên lưới 4$\times$4)};
  \scangrid{11.5}{hàng}{(0.5,0.5) (1.5,0.5) (2.5,0.5) (3.5,0.5) (0.5,1.5) (1.5,1.5) (2.5,1.5) (3.5,1.5) (0.5,2.5) (1.5,2.5) (2.5,2.5) (3.5,2.5) (0.5,3.5) (1.5,3.5) (2.5,3.5) (3.5,3.5)}
  \scangrid{14.5}{cột}{(0.5,0.5) (0.5,1.5) (0.5,2.5) (0.5,3.5) (1.5,0.5) (1.5,1.5) (1.5,2.5) (1.5,3.5) (2.5,0.5) (2.5,1.5) (2.5,2.5) (2.5,3.5) (3.5,0.5) (3.5,1.5) (3.5,2.5) (3.5,3.5)}
  \scangrid{17.5}{Morton}{(0.5,0.5) (1.5,0.5) (0.5,1.5) (1.5,1.5) (2.5,0.5) (3.5,0.5) (2.5,1.5) (3.5,1.5) (0.5,2.5) (1.5,2.5) (0.5,3.5) (1.5,3.5) (2.5,2.5) (3.5,2.5) (2.5,3.5) (3.5,3.5)}
  \node[note, anchor=north west, text width=8cm] at (11.5,-9.6)
    {Cây ghép từng cặp lá liên tiếp theo thứ tự quét. Morton gom khối 2$\times$2, 4$\times$4, 8$\times$8,\\nên mỗi qubit tầng trên ứng với một vùng vuông của ảnh.};
\end{tikzpicture}

\end{document}
